\makeabstract
    {
    Protease-cleavable masking strategy to reduce toxicity of de novo cytokine therapeutics 
    }
    {
    Robiah Arefin Ibn Mahmud\textsuperscript{1,2,3}, Julissa G. Tello\textsuperscript{3,4}, Frank Peprah\textsuperscript{3,4}, Hojeong Shin\textsuperscript{5}, Jung-Ho Chun\textsuperscript{5}, David Baker\textsuperscript{5}, Michael Dougan\textsuperscript{2,4}, Stephanie K. Dougan\textsuperscript{2,3}
    }
    {
    \textsuperscript{1} Department of Biology, Amherst College, Amherst, MA\\
\textsuperscript{2} Department of Immunology, Harvard Medical School, Boston, MA\\
\textsuperscript{3} Department of Cancer Immunology and Virology, Dana-Farber Cancer Institute, Boston, MA\\
\textsuperscript{4} Division of Gastroenterology, Department of Medicine, Massachusetts General Hospital, Boston, MA\\
\textsuperscript{5} Institute for Protein Design, Department of Biochemistry, University of Washington School of Medicine, Seattle, WA
    }
    {
    Cytokines are immune signaling proteins that can induce antitumor responses by stimulating cytotoxic lymphocytes (CTLs). However, clinical adoption of cytokine immunotherapies has been hindered by systemic immune activation and dose-limiting toxicity. We investigated a de novo protein-design strategy for synthesizing protease-activated cytokine prodrugs. These prodrugs consist of a therapeutic cytokine bound to a masking domain via a linker cleaved by tumor-associated matrix metalloproteinases (MMPs). In its masked state, the cytokine cannot associate with its cognate receptor and remains inactive. Once the linker is cleaved by MMPs, the cytokine regains pharmacological activity and stimulates tumor-targeting CTLs. We reasoned that conditional prodrug activation in tumors can limit systemic toxicity while retaining potent anti-tumor efficacy.
    }
    {
    We cultured a KPC cell line derived from murine pancreatic ductal adenocarcinoma (PDAC) and injected the cancer cells subcutaneously into 20 healthy wild-type mice. By day 13, the mice had developed palpable tumors and were tail-vein injected with 21h10 or Neo-2 treatments. Four treatment conditions were used for both cytokines: untreated (negative control), unmasked cytokine (positive control), masked cytokine (experimental group), and a masked construct with a non-cleavable linker (negative control). Three hours after injection, the mice were euthanized, and their spleens and tumors were harvested and lysed in RIPA lysis buffer. Western blot densiometric analysis was used to measure the ratio of phosphorylated STAT (P-STAT) to total STAT, quantified as fold change relative to untreated controls, using ImageJ and GraphPad Prism.
    }
    {
    In Neo-2-treated mice, masked Neo-2 induced negligible P-STAT5 signaling in the spleen but higher levels in the tumor, consistent with preferential tumor activation. In contrast, unmasked Neo-2 produced low P-STAT5 levels in both tissues, comparable to untreated controls. Similarly, masked 21h10 induced negligible P-STAT3 signaling in the spleen but higher levels in the tumor. Unmasked 21h10 produced low P-STAT3 levels in both tissues, comparable to untreated controls. 
    }
    {
    Our results suggest that masking enables preferential activation of cytokine prodrugs within tumors while limiting systemic signaling. Negligible STAT phosphorylation in the spleen following masked cytokine treatment, together with increased tumor signaling, supports the potential of this approach to reduce off-target immune activation and toxicity. However, further in vivo studies in tumor-bearing mice are needed to determine whether preferential tumor activation translates into improved therapeutic outcomes. Future experiments should assess anti-tumor efficacy through tumor growth and survival, and systemic toxicity through changes in body weight.
    }

    \newpage
    